% Part: first-order-logic
% Chapter: sequent-calculus
% Section: proof-theoretic-notions

\documentclass[../../../include/open-logic-section]{subfiles}

\begin{document}

\begin{editorial}
Cette section rassemble les définitions de la relation de
!!{derivability} et de la cohérence pour le calcul des séquents.
La mention de la déduction naturelle dans la source est une erreur
éditoriale corrigée ici.
\end{editorial}

\iftag{FOL}
      {\olfileid[fr]{fol}{seq}{ptn}}
      {\olfileid[fr]{pl}{seq}{ptn}}

\olsection{Notions de théorie de la démonstration}

\begin{explain}
Nous avons défini plusieurs notions sémantiques importantes : validité,
conséquence logique et satisfaisabilité. Définissons maintenant les
\emph{notions de théorie de la démonstration} correspondantes.
Elles font appel à la !!{derivability} ou à la !!{nonderivability}
de certains séquents, plutôt qu'à la satisfaction d'!!{sentence}s
dans des !!{structure}s. La coïncidence de ces notions constitue une
découverte importante : c'est le contenu des \emph{théorèmes de
correction et de complétude}.
\end{explain}

\begin{defn}[Théorèmes]
!!^a{sentence}~$!A$ est un \emph{théorème} s'il existe une
!!{derivation}, dans $\Log{LK}$, du séquent $\quad \Sequent !A$.
On écrit $\Proves !A$ si $!A$ est un théorème, et $\Proves/ !A$
dans le cas contraire.
\end{defn}

\begin{defn}[!!^{derivability}]
!!^a{sentence}~$!A$ est \emph{!!{derivable} à partir} d'un ensemble
d'!!{sentence}s~$\Gamma$, ce que l'on note $\Gamma \Proves !A$,
si et seulement s'il existe une partie finie $\Gamma_0 \subseteq \Gamma$
et une suite $\Gamma_0'$ énumérant les !!{sentence}s de $\Gamma_0$,
avec répétitions éventuelles, telles que $\Log{LK}$ !!{derive}s
$\Gamma_0' \Sequent !A$. Si $!A$ n'est pas !!{derivable} à partir
de $\Gamma$, on écrit $\Gamma \Proves/ !A$.
\end{defn}

Grâce aux règles de contraction, d'affaiblissement et d'échange,
l'ordre des !!{sentence}s et le nombre de leurs occurrences dans
$\Gamma_0'$ n'importent pas : si $\Gamma_0' \Sequent !A$ est
!!{derivable}, alors $\Gamma_0'' \Sequent !A$ l'est aussi pour
toute suite $\Gamma_0''$ contenant exactement les mêmes
!!{sentence}s que $\Gamma_0'$. Par exemple, si
$\Gamma_0 = \{!B, !C\}$, les suites
$\Gamma_0' = \tuple{!B, !B, !C}$ et
$\Gamma_0'' = \tuple{!C, !C, !B}$ contiennent toutes deux
exactement les !!{sentence}s de $\Gamma_0$. Si le séquent ayant
l'une pour antécédent est !!{derivable}, celui qui a l'autre pour
antécédent l'est aussi, comme le montre l'exemple suivant :
\begin{prooftree}
  \AxiomC{}
  \Deduce$!B, !B, !C \fCenter !A$
  \RightLabel{\LeftR{\Contraction}}
  \UnaryInf$!B, !C \fCenter !A$
  \RightLabel{\LeftR{\Exchange}}
  \UnaryInf$!C, !B \fCenter !A$
  \RightLabel{\LeftR{\Weakening}}
  \UnaryInf$!C, !C, !B \fCenter !A$
\end{prooftree}
Désormais, si $\Gamma_0$ est un ensemble fini d'!!{sentence}s,
nous noterons $\Gamma_0 \Sequent !A$ un séquent dont l'antécédent
est une suite énumérant les éléments de $\Gamma_0$, avec répétitions
éventuelles. Nous laisserons implicites les contractions, échanges
et affaiblissements nécessaires.

\begin{defn}[Cohérence]
Un ensemble d'énoncés~$\Gamma$ est \emph{incohérent} s'il existe
une partie finie $\Gamma_0 \subseteq \Gamma$ telle que $\Log{LK}$
!!{derive}s $\Gamma_0 \Sequent \quad$. S'il n'est pas incohérent,
c'est-à-dire si, pour toute partie finie $\Gamma_0 \subseteq \Gamma$,
$\Log{LK}$ ne permet pas de !!{derive} $\Gamma_0 \Sequent \quad$,
on dit que $\Gamma$ est \emph{cohérent}.
\end{defn}

\begin{prop}[Réflexivité]
\ollabel{prop:reflexivity}
Si $!A \in \Gamma$, alors $\Gamma \Proves !A$.
\end{prop}

\begin{proof}
Le séquent initial $!A \Sequent !A$ est !!{derivable}, et
$\{!A\} \subseteq \Gamma$.
\end{proof}

\begin{prop}[Monotonie]
\ollabel{prop:monotonicity}
Si $\Gamma \subseteq \Delta$ et $\Gamma \Proves !A$, alors
$\Delta \Proves !A$.
\end{prop}

\begin{proof}
Supposons $\Gamma \Proves !A$. Il existe donc une partie finie
$\Gamma_0 \subseteq \Gamma$ telle que $\Gamma_0 \Sequent !A$
soit !!{derivable}. Comme $\Gamma \subseteq \Delta$,
$\Gamma_0$ est aussi une partie finie de $\Delta$. La
!!{derivation} de $\Gamma_0 \Sequent !A$ montre donc également
que $\Delta \Proves !A$.
\end{proof}

\begin{prop}[Transitivité]
\ollabel{prop:transitivity}
Si $\Gamma \Proves !A$ et $\{!A\} \cup \Delta \Proves !B$,
alors $\Gamma \cup \Delta \Proves !B$.
\end{prop}

\begin{proof}
Si $\Gamma \Proves !A$, il existe une partie finie
$\Gamma_0 \subseteq \Gamma$ et une !!{derivation}~$\pi_0$
de $\Gamma_0 \Sequent !A$. Si $\{!A\} \cup \Delta \Proves !B$,
il existe une partie finie $\Delta_0 \subseteq \Delta$ et une
!!{derivation}~$\pi_1$ de $!A, \Delta_0 \Sequent !B$,
en ajoutant au besoin $!A$ par affaiblissement. Considérons alors
la !!{derivation} suivante :
\begin{prooftree}
  \AxiomC{}
  \RightLabel{$\pi_0$}
  \Deduce$\Gamma_0 \fCenter !A$
  \AxiomC{}
  \RightLabel{$\pi_1$}
  \Deduce$!A, \Delta_0 \fCenter !B$
  \RightLabel{\Cut}
  \BinaryInf$\Gamma_0, \Delta_0 \fCenter !B$
\end{prooftree}
Comme $\Gamma_0 \cup \Delta_0 \subseteq \Gamma \cup \Delta$,
on obtient $\Gamma \cup \Delta \Proves !B$.
\end{proof}

En particulier, si $\Gamma \Proves !A$ et $!A \Proves !B$, alors
$\Gamma \Proves !B$. Il en résulte aussi que si
$!A_1, \dots, !A_n \Proves !B$ et $\Gamma \Proves !A_i$
pour chaque~$i$, alors $\Gamma \Proves !B$.

\begin{prop}
\ollabel{prop:incons}
$\Gamma$ est incohérent si et seulement si $\Gamma \Proves {!A}$
pour tout énoncé~$!A$.
\end{prop}

\begin{proof}
La démonstration est laissée en exercice.
\end{proof}

\begin{prob}
Démontrer la \olref{prop:incons}.\footnote{Note de l'édition
française : le renvoi est relatif à la section courante, afin de
fonctionner dans les présentations propositionnelle et du premier ordre.}
\end{prob}

\begin{prop}[Compacité]
  \ollabel{prop:proves-compact}
  \begin{enumerate}
  \item Si $\Gamma \Proves !A$, il existe une partie finie
    $\Gamma_0 \subseteq \Gamma$ telle que $\Gamma_0 \Proves !A$.
  \item Si toute partie finie de $\Gamma$ est cohérente,
    alors $\Gamma$ est cohérent.
  \end{enumerate}
\end{prop}

\begin{proof}
  \begin{enumerate}
  \item Si $\Gamma \Proves !A$, il existe une partie finie
    $\Gamma_0 \subseteq \Gamma$ telle que le séquent
    $\Gamma_0 \Sequent !A$ possède une !!{derivation}.
    Par conséquent, $\Gamma_0 \Proves !A$.
  \item Si $\Gamma$ est incohérent, il existe une partie finie
    $\Gamma_0 \subseteq \Gamma$ telle que $\Log{LK}$
    !!{derive}s $\Gamma_0 \Sequent \quad$. Ainsi, $\Gamma_0$
    est une partie finie de $\Gamma$ qui est incohérente.
  \end{enumerate}
\end{proof}

\end{document}
